% discussion.tex -- Findings written 2026-10-29 for the label NEG of P (SAP Addendum A4, A4.2.3: opening of 6.1 near-verbatim).
% The fusion-context and limitations subsections are outcome-independent and drafted from the preregistration (card v4, SAP, A4).
% No sentence in this file draws on PS, S2, B, R, E-dec, E-pred, the suite, detection or the exploratory comparison to explain, qualify
% or predict the primary result (AB.1). The phrase "when does input gating help" is not used (D18 reading 1, decisions D0076-D0077);
% after the opening, the question of when gating can help is discussed only through the closed-form condition and the P table.
% Proposition 2 (correlated tokens) is not used on E-pred (D0078); Proposition 1 of the theory addendum is not used (D0078-D0079).

\subsection{Findings}\label{ssec:disc_findings}
% run: P2-GATE-M4 (analysis.json primary.*)
The preregistered test excludes an improvement or a loss of more than five percentage points in image-direction verification accuracy
from gating under JPEG quality~10, Gaussian noise with $\sigma=0.10$ and a $128$-pixel splice, for this model and these severities; it
does not show that gating has no effect.

The primary table shows how narrow the measured change is: in 500 items the gate changed the decision once under JPEG, never under
noise and four times under splice (Table~\ref{tab:res_primary}).
The exclusion of $\pm0.05$ is a statement about these discordances at the threshold calibrated for $\alpha=0.05$, at one severity per
corruption, in one model and on one data set.
It says nothing about a gate that flags a larger share of tokens, about stronger corruptions, or about another model or task, none of
which was tested.
Section~\ref{ssec:evidence} gives, under independent tokens, the condition on the hit and false-alarm rates under which discarding
flagged tokens raises the signal-to-noise ratio of an additive evidence sum (Eq.~\eqref{eq:snr}).
The condition states what a gate would have to achieve; it is an approximation for dependent image tokens, the test was not designed to
check it, and we do not use it to account for the result.

The remaining findings are separate measurements, each under its own family label, and none of them is an account of the primary result.
For each we give the observation, a hypothesis for it, and what would falsify the hypothesis.

\emph{The clean false-alarm statement held; the flags on corrupted images are not covered by it.}
% run: P2-GATE-M4 (calibration.*)
On held-out clean pairs the realised flagged share was 0.0456 at $\alpha=0.05$ (95\,\% interval 0.0442 to 0.0470), with no evidence
of violation at any of the three levels (Section~\ref{ssec:res_cal}).
On the unchanged tokens of corrupted images the false-alarm proportion was 0.036 to 0.044 for the realistic corruptions and 0.089 under
swap, above $\alpha$.
The conformal statement covers clean pairs only, so these proportions lie outside it; the audit of the positive control
did not treat the 0.089 as a defect for that reason.
Hypothesis: a token next to a replaced code is scored with damaged context and is flagged more often, as Section~\ref{ssec:calibration}
anticipates.
It would be falsified if, under swap, the false-alarm proportion of an unchanged token did not fall with its distance to the nearest
replaced code; this was not measured.

\emph{The gate's score was near chance at finding the codes that JPEG and noise change.}
% run: P2-GATE-M1; P2-GATE-M4 (s_d2.*, calibration.per_condition.*)
On the detection items and again on the test items (families S-D and S-D2) the per-item AUROC of $q$ was 0.561 under JPEG and 0.503 and
0.509 under noise, below a texture score whose sign was chosen on the data; under noise the hit rate of the gate (0.033) was below its
false-alarm proportion on unchanged tokens (0.036).
Hypothesis, the label-circularity concern of Section~\ref{ssec:limits}: JPEG quality~10 and noise change more than nine in ten codes of
an image (Table~\ref{tab:res_flags}), so a changed code marks where the tokenizer is unstable more than where the image is damaged, and a
texture score tracks that instability.
It would be falsified if $q$ stayed at chance against a damage label defined in the pixel domain, for example the reconstruction error
of each token's patch, under which texture had no advantage; this was not measured.

\emph{A pixel denoiser restored most of the accuracy lost to noise.}
% run: P2-GATE-M4 (s_dn.*, c1_test.accuracy)
Under noise the denoised copy reached 0.852 against 0.676 without a gate and 0.880 on clean images, and it gave higher verification
accuracy than GATE-M (family S-DN, Section~\ref{ssec:res_dn}); under JPEG no difference was detected, with a denoiser setting that
lowers JPEG image quality.
Hypothesis: noise of this strength is largely removed in the pixel domain before tokenisation, so the tokens of the denoised image are
close to the clean ones.
It would be falsified if the share of codes that differ from the clean tokenisation were not lower for the denoised noisy images than for
the noisy ones; this was not measured.
The denoiser is a measured comparator, not part of the gate, and one setting was used.

\emph{Withholding every changed code lowered accuracy under JPEG.}
% run: P2-GATE-M4 (mechanism.jpeg10.y.changed_code_minus_none; edec.json cells.P1; calibration.per_condition.jpeg10.prop4.rho)
The changed-code gate, which keeps only the codes that the corruption left unchanged, lowered accuracy under JPEG by 0.084 and raised it
under noise by 0.070 (Table~\ref{tab:res_mech}); in E-dec the gap to this oracle was negative under JPEG with an interval that excludes
zero, which by the plan voids any detection reading of that cell (Section~\ref{ssec:res_edec}).
Hypothesis: under JPEG this gate decides on about four percent of the tokens of an image, and the sum over so few tokens is noisier than
the sum over all of them, which is the composition effect named in the plan.
It would be falsified if keeping a random set of the same small size, instead of the unchanged codes, did not lower accuracy by a similar
amount; this was not measured.

\emph{Exploratory: no detector-free combiner was above the hard gate on a realistic corruption.}
% run: P2-GATE-M4 (tier3_fusion.json)
In the exploratory comparison (Section~\ref{ssec:res_expl}; Tier~3, intervals only), the hard gate gave higher verification accuracy
than the median and the 20\,\% trimmed mean under JPEG and splice; the data do not show a difference between the hard gate and the
trimmed mean at the gate's flag count on any realistic corruption, nor between the hard gate and the median or the 20\,\% trimmed mean
under noise, with bounds at most 0.05, which counts by the pre-stated rule as a finding against the added value of the detector on those
cells.
Discarding the flagged tokens and discounting or weighting them were not distinguishable on any realistic corruption.
This is a property of these combiners on these cells, measured on the same evidence as the primary test, and not an account of the
primary result; when the gate's score is near chance, the comparison with the trimmed mean at the gate's flag count largely repeats the
primary contrast (limit (c) of Table~\ref{tab:res_tier3}).

\emph{The positive control was not met; the audit it triggered found no defect.}
% run: P2-GATE-M4 (ps.*); P2-A23-AUDIT
On the two sanity corruptions the gate did not change accuracy detectably (Section~\ref{ssec:res_ps}): positive control not met; audit
triggered (SAP A2.3); no defect found, with $n_d=4$ (patch) and $n_d=13$ (swap) discordant items.
The audit covers the classes of error that the plan names and does not validate the pipeline; Section~\ref{ssec:res_ps} states what
it cannot see.
Nothing in the positive control is used in this paper to account for the primary result.

\emph{Post hoc and descriptive: the decision headroom of the positive control.}
% run: P2-A23-AUDIT (a23_ps_audit.json, per_condition.{patch,swap}); P2-GATE-M4 (ps.*)
This paragraph was written after the outcome of PS was known, enters no test, and describes the task and the design of PS, not the
gate.
Without a gate, the sanity corruptions changed the decision on 5 of the 500 test items under patch (all five from right to wrong) and
on 9 under swap (7 from right to wrong, 2 from wrong to right), that is on 0.010 and 0.018 of the items.
With no losses, the exact two-sided McNemar test at the first Holm step (level 0.025) needs at least 7 wins; a gate that returned every
item to its decision on the clean image would have given $b=5$, $c=0$, $p=0.0625$ under patch and $b=7$, $c=2$, $p=0.18$ under swap,
and would not have reached that step.
The changed-code arm, which uses the true set of replaced codes, did not either ($b=4$, $c=2$, $p=0.69$ under patch; $b=7$, $c=4$,
$p=0.55$ under swap), and GATE-M made the same decision as that arm on 498 of the 500 items in each cell.
The discordance between GATE-M and no gate was 0.008 under patch and 0.026 under swap, against the 0.05 to 0.20 that the power
calculation of the plan assumed (AD.4, where it is marked as an assumption).
At that discordance this test could not have rejected for a perfect gate, so the positive control, as designed, cannot separate a
working pipeline from a broken one on these two corruptions; this is a property of the task, where replacing 64 or 154 codes rarely
turns a correct decision into a wrong one, and of the design, which assumed more discordance than occurred, and not a statement about
the method.
The description is not used to account for the primary result, and the positive control is not read as met on this ground.

\subsection{Where the work sits in information fusion}\label{ssec:disc_fit}
\emph{What is fused.}
The 1\,024 image tokens of one image are 1\,024 sources.
Each reports a log-likelihood ratio between two candidate captions, and the decision sums them; a second, text-direction statistic can
be fused with that sum at score level (family B).
\emph{At what level.}
At the input: the gate acts on the sources before the evidence sum, not on the output decision, which is why the paper speaks of
input-level source validation.
\emph{Where reliability enters.}
Only through a zero-one weight in the sum (GATE-M) or through replacement of the flagged values by a model draw (GATE-I).
The weight is computed from within-modality evidence alone, and its threshold is set on clean data at a fixed false-alarm level, as in
CFAR detection.
\emph{What is not claimed.}
The paper claims no new fusion rule.
The additive rule is optimal only for sources that are independent given the hypothesis, which image tokens are not, and the hard
weight is not claimed to be the best way to use a reliability estimate; graded weights and robust aggregation are the standard
alternatives, and Section~\ref{ssec:tier3} states the exploratory comparison with them.

The calibrated gate of Zhou and Lu~\cite{zhou2027burden} and the gate studied here answer different questions.
As described in the abstract of Zhou and Lu, theirs decides which of two recognisers to believe for one test instance; the gate
here decides which input tokens to withhold before two directions of image-text evidence are combined.
The two could be stacked, but this is not tested.
Evidential fusion~\cite{shao2024ddef,huang2025evidential} and conformal prediction sets on a fusion network~\cite{wang2025multistage}
are other uses of reliability and of distribution-free calibration in this journal that we found; as far as their abstracts describe,
none scores input tokens with the likelihood of a generative model.
\todo{AUTHOR: after the results are frozen, give the abstract and Figure~1 caption to a reader with no project context and ask what is
fused (open check on the scope-fit question); settle the scope-fit question with the supervising professor before submission.}

\subsection{Limitations}\label{ssec:limits}
Except for the two items on the positive control, the limitations below are fixed by the design and do not depend on the outcome.
\begin{itemize}
\item \emph{One model, one data set.} All measurements of the primary family use UniDisc 1.4B and COCO val2017 images as tokenised
by one tokenizer. A second, 4-bit quantised model was to enter only in family R, only if P had a rejection with $b>c$; P had none, so R
was not run. No statement is made for other models, resolutions, caption styles or data sets.
\item \emph{Contamination is unknown.} Whether the evaluation images were in the training data of the model or of the tokenizer was not
established, and it could make clean pairs easier than unseen ones. The same applies, with more force, to the hard-negative set,
whose true captions are COCO captions.
\item \emph{Narrow guarantee.} The calibration bounds the expected false-alarm proportion on clean, exchangeable test pairs. It gives
no bound on detection power, none on the clean tokens of corrupted items, and none under shift. Unlike a selective-risk bound on
the harm of an action, as described for Zhou and Lu~\cite{zhou2027burden}, it says nothing about the harm or the power of the gate
on corrupted inputs. The realised proportion on one finite test set can lie above $\alpha$ because the statement is marginal over the
calibration draw.
\item \emph{Corruptions and labels.} Corruption labels are relative to the tokenisation of the clean image, so a ``clean'' token
may already carry compression artefacts of the original photograph. Changed codes under JPEG and noise concentrate where the tokenizer
is unstable, so the label may track tokenizer instability more than damage: a texture rival that wins against this label is not
thereby shown to explain the corruption effect, and the changed-code gate is not an upper bound for these corruptions.
\item \emph{Comparators chosen on the data.} The best cheap rival is selected among about nine candidates, with both signs admitted, which
biases the detection contrast against the gate. The rival's sign could not be fixed blind (Section~\ref{ssec:prereg}).
\item \emph{Pixel denoiser.} One setting was used for all corruptions; it was chosen after one exploratory evaluation on 20 images and
harms JPEG images, so the denoising contrast on JPEG favours the gate.
\item \emph{Variance not in the intervals.} Intervals cover item sampling. They do not cover variation across mask partitions, corruption
draws, jitter seeds or models, because no training seeds exist and only one corruption draw is made per image. Eleven test families
are run with no multiplicity control across them.
\item \emph{A constructed task.} The primary decision is a two-way caption match on the weaker of the two available statistics, with
random negatives (Section~\ref{ssec:task}). It measures how an image-direction statistic tolerates damage and says nothing about
what a gate adds to a verifier that uses the stronger text direction; only family B addresses the two-direction setting, and it may be
untestable at the ceiling or not run.
\item \emph{A hard gate and dependent tokens.} The gate is binary, and the evidence sum and the signal-to-noise analysis assume
independent tokens, which neighbouring image tokens are not. Graded and robust alternatives are examined only in an exploratory
comparison (Section~\ref{ssec:tier3}), which is not a test and is subject to its own limits.
\item \emph{Time input and checkpoint.} The scorer calls the denoiser without a time input; that this is the intended path was checked
for the released configuration and checkpoint only, the training weight of whole-modality masks is unknown, and the key check does not
cover the image tokeniser (Section~\ref{ssec:data}, \ref{app:timefree}).
\item \emph{Scope of the preregistered run.} Families S2 (re-infilling), B (two-way fusion) and R (second model) each have a status of
run or not run, decided by time and by their gatekeeping rule and never by an outcome, and the paper reports that status.
Caption corruption, the backbone authors' own re-infill recipe and the gating of caption tokens are not tested as families.
\item \emph{Positive control not met.} The preregistered positive control on the sanity corruptions was not met, which triggered a
pipeline audit (Section~\ref{ssec:res_ps}); the audit found no defect. It was run after the outcome was known and covers the classes
of error that the plan names. A recomputation from the stored arrays cannot see an error upstream of them, and the audit does not
validate the pipeline. The conformal statement covers clean pairs only; on the unchanged tokens of swapped images the flag proportion
was 0.089, above $\alpha$, which the audit recorded and did not count as a defect.
\item \emph{Power assumption of the positive control.} The plan's power calculation for the positive control assumed a discordance of
0.05 to 0.20 between the gated and the ungated decision. The realised values were 0.008 and 0.026 ($n_d=4$ and 13 of 500), and at
that headroom the preregistered test could not have rejected even for a gate that restored every clean decision. The positive control
therefore cannot separate a working pipeline from a broken one on these two corruptions. This is a property of the design and of the
task, not of the method, and it is not used to account for the primary result.
\item \emph{Hardware and missing data.} The GPU is a single laptop device that failed in earlier pilots. A crash leaves a prefix of the item
list and the missing-item rule applies; a fold under 240 items is reported as incomplete.
\item \emph{Narrow novelty.} Detecting and remasking corrupted input tokens, cross-modal remasking, risk-controlled thresholds for
vision-language errors, including gates that decide between whole predictions~\cite{zhou2027burden}, and masked-likelihood scores
each have precedents (Section~\ref{sec:related}). The claim is a measurement of their composition on one model, and it holds only as
far as the measured results go.
\item \emph{Concurrent work.} A search for concurrent work was last run on 28 October 2026; unpublished or under-review work is not
visible, and the closest paper in this journal~\cite{zhou2027burden} was compared from its abstract and preview only.
\todo{AUTHOR: repeat the search once before submission, and read the full text of Zhou and Lu.}
\end{itemize}
